% \subsection{Overview}\label{block: overview}
% \begin{table*}[t!]
%     % \footnotesize
%     \centering
%     \caption{Summary of operation costs for homomorphic SQL operators.
%     Here, the $\dbtabrow$ stands for the number of rows in the table 
%     $\dbtab$, the $\abs{\attribute}$ represents the range of 
%     attribute.}
%     \vspace{-0.5em}
%         \begin{tabular}{ccccc}
%             \toprule
%             operator & \#\homcomp & \#\homgate & \#\cmux & \#\cb\\
%             \midrule
%             $\select$ & $\dbtabrow \cdot \dbabsquery$& 
%             $\abs{\dbtab}_r \cdot (\dbabsquery - 1)$ & 0 & 0\\
%             $\join$ & $\abs{\dbtab_a}_c \cdot \abs{\dbtab_b}_c$&
%             0 & 0 & 0\\
%             $\groupby$ & $\dbtabrow \cdot \abs{\attribute}$ &
%             0 & 0 & 0\\
%             $\orderby$ & $\dbtabrow \log_{2}^{2}{\dbtabrow } $ & 0 
%             & $\dbtabrow (\log_{2}^{2}{\dbtabrow } + 1) $&
%             $\dbtabrow (\log_{2}^{2}{\dbtabrow } + 1) $ \\
%              $\dbmin/\dbmax$ & $\dbtabrow - 1$ & 0 & $2\dbtabrow - 1$ & $2\dbtabrow - 1$\\
%             
%             
%             \bottomrule
%         \end{tabular}
%     \label{tab: sqlcost}
% \end{table*}
% 
% In this section, we talk about how \NAME supports SQL query
% operators including $\select, \join, \groupby, \orderby$
% , and aggregation functions. 
% We provide a simple plaintext illustration in \cref{fig: query_eva}
% and analyze the homomorphic operator costs for distinct SQL statements with respect to the input size $\dbabstab$ in \cref{tab: sqlcost}.
% 
% $\bullet$ $\boldsymbol{\select}$. 
% The homomorphic $\select$ operator inputs 
% $\dbabsquery$ encrypted predicates with 
% $\dbtabcolumn$ columns of RLWE ciphertexts
% and outputs a set of LWE ciphertexts
% $\mathbf{sc} = (sc_0, sc_1 ,..., sc_{\dbtabrow-1})$
% where each $sc_i$ is an LWE
% ciphertext which encrypts $1$ or $0$, 
% indicating the $i$-th row of $\dbtab$
% is selected or not.
% The homomorphic $\select$ operator contains 
% the homomorphic predicate evaluation and the gate operation
% between the predicate results.
% In \NAME, 
% the encrypted predicate (as shown in \cref{fig: query_eva})
% is the logical comparison 
% ($>, \geq, <, \leq, ==, \ne$)
% between 
% two ciphertexts
% and could be evaluated by the
% homomorphic comparison $\homcomp$ operator in \cref{block: comparison}.
% The gate operation (AND, OR)
% between the comparison results
% could be implemented by the $\homgate$ defined in 
% \cref{pre:he_circuits}.
% The cost of $\homcomp$ in $\select$ is $\dbtabrow \cdot 
% \dbabsquery$ since 
% each predicate requires the comparison between all the encrypted records ($\dbtabrow$) of 
% the certain column and 
% the ciphertext of the predicate value.
% And the cost of $\homgate$ in $\select$ is $\dbtabrow \cdot (\dbabsquery - 1)$  
% as the $\dbabsquery$ predicates contains $(\dbabsquery - 1)$
% gate operators applied to the results of the comparison.
% % Consider an example 
% % predicate \texttt{ID == 2158} in $\dbquery$,
% % it requires the comparison between all the encrypted records ($\abs{\dbtab}_r$) of 
% % the column \texttt{ID} and 
% % the ciphertext of the predicate value \texttt{2158},
% % which contains $\abs{\dbtab}_r$ homomorphic comparisons.
% % Moreover, the results of $\abs{\dbquery}$ predicates
% % will be integrated into the selected result by
% % $\abs{\dbquery} - 1$ gate operators, while each operators handles
% % $\abs{\dbtab}_r$ comparison results from homomorphic predicate evaluation.
% % in \cref{block: comparison}. 
% 
% $\bullet$ $\boldsymbol{\join}$.
% A homomorphic inner join is composed of 
% the Cartesian product of the two encrypted tables
% and the evaluation of the join condition, which could be 
% implemented by homomorphic comparison operator $\homcomp$.
% Thus, the cost of  $\homcomp$ in $\join$ is 
% $\abs{\dbtab_a}_c \cdot \abs{\dbtab_b}_c$.
% 
% $\bullet$ $\boldsymbol{\groupby}$.
% To evaluate the homomorphic $\groupby$ operator,
% we segregate the original query $\dbquery$ into several queries, 
% each with additional equality comparison predicates.
% For example, $\groupby$ the attribute $\mathsf{Sex = \{0,1\}}$ 
% could be implemented by $\dbquery$ with additional predicate $\mathsf{Sex = 0}$
% and $\dbquery$ with additional predicate predicate $\mathsf{Sex = 1}$.
% Supposing the range of $\groupby$ attribute $\attribute$ is $\abs{\attribute}$.
% Thus, the cost of  $\homcomp$ in $\groupby$ is 
% $\dbtabrow \cdot \abs{\attribute}$.
% 
% $\bullet$ $\boldsymbol{\orderby}$.
% The homomorphic $\orderby$ operator could be implemented by
% homomorphic comparison $\homcomp$ and homomorphic selector $\cmux$.
% The main challenge is how to only sort the selected encrypted records. 
% We detail our method to overcome this challenge \cref{block: minmax}.
% 
% $\bullet$ \textbf{Aggregation functions}.
% As shown in \cref{fig: query_eva},
% after the homomorphic predicate evaluation, 
% we get $\abs{\mathcal{T}}_{r}$
% LWE ciphertexts encrypt the filtering results.
% The homomorphic aggregation computes the functions
% on the filtered rows.
% \NAME supports both logic aggregation 
% (such as $\dbmin, \dbmax$, and $\dbtopk$)
% and arithmetic aggregation functions
% (such as $\dbsum, \dbavg$, and $\dbcount$).
% We detail the exact method in \cref{block: minmax} and \cref{block: LWEsPack}.
% As shown in \cref{fig: query_eva}, to calculate \texttt{MIN(HR)}, 
% we first selected the satisfied rows and perform \texttt{HomMIN} operations.}

% \revise{\cref{fig: query_eva} illustrate a plaintext process 
% of query evaluation in \NAME.
% It works on the encrypted table $\mathcal{CT}$, that is, 
% a set of RLWE ciphertexts 
% $(ct_1, ct_2, ..., ct_{\lvert \mathcal{T} \rvert_{c}})$.
% On a high level, the evaluation of encrypted query 
% $\mathcal{CQ} = (cp_{1}, cp_{2}, ... ,cp_{\lvert \mathcal{Q} \rvert}, B, G)$ 
% on $\mathcal{CT}$  
% can be summarized in the following steps:}

% \revise{\textbf{1) homomorphic predicate evaluation}: 
% For each encrypted predicate 
% $cp_i \in \mathcal{CQ}, i \in (1, 2, ..., \abs{\mathcal{Q}})$,
% we take the matched encrypted column ciphertext 
% $ct_d \in \mathcal{CT}$ and the predicate value ciphertext $cp_i$, 
% which are two RLWE ciphertexts, 
% to perform \texttt{HomCom} operations 
% and get a set of LWE ciphertexts $\mathbf{ctp_i}$ 
% where the $j$-th LWE ciphertext of $\mathbf{ctp_i}$ indicates 
% the $j$-th row of the table is satisfied the predicate $cp_i$ or not.
% Then we calculate the \texttt{HomGate} 
% between each comparison result $\mathbf{ctp_i}$.
% After the evaluation, we obtain a set of LWE ciphertexts 
% $\mathbf{fr} = (fr_1, fr_2 ,..., fr_{\abs{\mathcal{T}}_r})$
% where each $fr_i$ is an LWE
% ciphertext which encrypts $1$ or $0$, 
% indicating the $i$-th row of $\mathcal{T}$
% is selected or not.}
% $\mathbf{}$
% to perform
% The homomorphic predicate evaluation of $\mathcal{CQ}$
% consists of the homomorphic comparison \texttt{HomComp} in each
% predicate $cp_i$ and the homomorphic gate \texttt{HomGate} 
% between predicates.
% \revise{Gate Tree..., add an algorithm, and complexity analysis...}
 
% $i \in (1, 2, ..., \abs{\mathcal{Q}})$ in $\mathcal{CQ}$, 
% we first get the attribute index of the predicate $d$ 
% and the comparison operator $op$
% ($\geq, \leq, ==, >, <$) of the predicate.
% Here, the $ct_d$ represents the RLWE ciphertext which
% encrypts the matching column in $\mathcal{T}$
% and the $cp_i$ refers to the RLWE ciphertext of the predicate value.
% Then, we sample extract the RLWE ciphertexts $ct_d$ and $cp_i$, 
% and get two LWE ciphertext vectors
% $\mathbf{ct_d}$ and $\mathbf{cp_i}$ wich size $\abs{\mathcal{T}}_{r}$.
% For each LWE ciphertext $ct_d[j]$and $cp_i[j]$, 
% we perform \texttt{HomCom} of the two LWE ciphertext
% with the comparison operator $op$
% and get the encrypted comparison result $comp_{i,j}$.
% Next, we evaluate the logic gate between different predicates 
% using \texttt{HomGate}.


% \revise{\textbf{2) homomorphic logic aggregation}:
% After the homomorphic predicate evaluation, 
% we get $\abs{\mathcal{T}}_{r}$
% LWE ciphertexts encrypt the filtering results.
% The homomorphic logic aggregation computes the logic function
% (such as \texttt{MIN/MAX}, \texttt{Top-k})
% on the filtered rows.
% As shown in \cref{fig: query_eva}, to calculate \texttt{MIN(HR)}, 
% we first selected the satisfied rows and perform \texttt{HomMIN} operations.}
% % the \texttt{HomMIN}
% % The $\texttt{LWEsPack}$
% % packs $\abs{\mathcal{T}}_{r}$ into a RLWE ciphertext $fl$
% % for latter homomorphic aggregation.

% \revise{\textbf{3) homomorphic arithmetic aggregation}:
% The aggregation process computes the arithmetic function
% (such as \texttt{SUM}, \texttt{AVG}, and \texttt{COUNT})
% on the filtered rows.
% As the filtered result format (LWE ciphertexts) is inefficient
% to homomorphic arithmetic function.
% We need to perform homomorphic ciphertext conversion 
% \texttt{HomConvert} to pack LWE ciphertexts into an RLWE ciphertext
% for latter operators.}
% For each aggregation filed $g_i$, 
% we get the attribute index $d = g(a)_i$
% and the aggregation function $f_i = g(f)_i$.
% As the filtering result $fl$ is in coefficient format, 
% we directly perform the aggregate function when $f_i$ supports 
% coefficient format computation.
% Otherwise, we first evaluate the \texttt{CoeffToSlot} 
% to convert $fl$ to $\hat{fl}$ slot format ciphertext $\hat{fl}$
% and then computes the $f_i$ in slot format.



